%% kapa: exam=2 round=32 year=2021 period=1 subject=appraisal_practice subject_ko=감정평가실무 seq=2 number=2 points=30
%% source: Q-Net 감정평가사 기출문제 (cst00309) artlSeq=5208993 file=raw/1008/5208993/감평 2차 시험지(2021).pdf
%% conversion: kordoc 4.24.51 -> tools/gen2.py (auto); review=proofread-sonnet; flags=none
\begin{longquestion}{【문제 2】}
소송감정인인 감정평가사 甲은 부당이득반환청구와 관련된 소송에서 토지에 대한 임대료의 감정평가를 의뢰받았다. 본 사건에서 토지 임대료에 대한 감정평가는 이미 다른 감정인에 의하여 완료되어 해당 재판부에 제출된 상황인데, 감정평가 결과에 대하여 피고는 부당함을 주장하였고 이것이 받아들여져 재의뢰가 된 사안이다. 제시된 자료를 활용하여 각 물음에 답하시오. (30점) \par
\medskip\noindent 물음 1) 감정평가 관련 법령 및 이론에 비추어 피고가 제기한 주장의 타당성 여부 및 그 근거를 약술하시오. (5점) \par
\medskip\noindent 물음 2) < 자료 3 > 실지조사, 자료수집 및 검토내용에 따라 시장가치에 기초한 기초가격에 적용할 기대이율(필요제경비 불포함)을 산출하시오. (15점) \par
\medskip\noindent 물음 3) 위 물음 2)에서 산출된 기대이율(시장가치 기준)을 기초로 본건의 연도별 적산임료를 구하시오. (5점) \par
\medskip\noindent 물음 4) 적산법의 장·단점 및 적용상 유의사항에 대하여 약술하시오. (5점) \par
\medskip\noindent < 자료 1 > 종전 감정평가 내역(요약) \par
\medskip\noindent 1. 감정사항: 경기도 K시 H동 104-2 토지(전) 1,652m\textsuperscript{2}에 대한 2018년 5월 1일부터 2021년 4월 30일까지의 임료 \par
2. 대상물건의 개요 \par
\hspace*{1em}가. 소재지: 경기도 K시 H동 104-2 \par
\hspace*{1em}나. 지목 및 이용상황: 전 / 전 \par
\hspace*{1em}다. 면적: 1,652m\textsuperscript{2} \par
\hspace*{1em}라. 토지이용계획사항: 자연녹지지역, 개발제한구역 \par
\hspace*{1em}마. 인근환경: 본건 토지는 서울특별시 북서측에 소재하는 경기도 K시 도심 남동측에 위치하는 근교농경지대에 소재하며, 부근 일대는 대부분 개발제한구역으로 원예농업을 위한 농경지로 이용중임 \par
3. 감정평가액의 산출 \par
\hspace*{1em}가. 감정평가방식의 적용: 토지에 대한 임대사례가 희박하여 적산법을 적용하되 다른 방식에 의한 검토는 생략함 \par
\hspace*{1em}나. 기초가격(공시지가기준법 적용) \par
{\small\centering
\begin{tabular}{|c|c|c|c|}\hline
기간 & 단가(원/m\textsuperscript{2}) & 면적(m\textsuperscript{2}) & 기초가격(원) \\\hline
2018. 5. 1.\textasciitilde{}2019. 4. 30. & 900,000 & 1,652 & 1,486,800,000 \\\hline
2019. 5. 1.\textasciitilde{}2020. 4. 30. & 956,000 & 1,652 & 1,579,312,000 \\\hline
2020. 5. 1.\textasciitilde{}2021. 4. 30. & 1,016,000 & 1,652 & 1,678,432,000 \\\hline
\end{tabular}\par}
\hspace*{1em}다. 기대이율: 1.0\%/연(기대이율 적용기준율표 참작) \par
\hspace*{2em}* 기대이율 적용기준율표(일부 발췌) - 감정평가 실무매뉴얼(임대료 감정평가편) \par
{\small\centering
\begin{tabular}{|c|c|c|c|}\hline
\multicolumn{2}{|c|}{대분류} & 소분류 & 실제이용상황 \\\hline
\multirow{2}{*}{Ⅱ} & \multirow{2}{*}{농지} & 도시근교농지 & 1.0\% 이내 \\\cline{3-4}
 &  & 기타농지 & 1.0\%\textasciitilde{}3.0\% \\\hline
\end{tabular}\par}
\hspace*{1em}라. 필요제경비: 보유세(재산세 등) 금액이 미미하고 기타 필요제경비가 필요하지 아니한 것으로 판단하여 기대이율에 포함하였음 \par
\hspace*{1em}마. 감정평가액(보증금 없는 상태 기준) \par
{\small\centering
\begin{tabular}{|c|c|c|c|}\hline
기간 & 기초가격(원) & 기대이율(\%) & 적산임료(원) \\\hline
2018. 5. 1.\textasciitilde{}2019. 4. 30. & 1,486,800,000 & 1.0 & 14,868,000 \\\hline
2019. 5. 1.\textasciitilde{}2020. 4. 30. & 1,579,312,000 & 1.0 & 15,793,120 \\\hline
2020. 5. 1.\textasciitilde{}2021. 4. 30. & 1,678,432,000 & 1.0 & 16,784,320 \\\hline
합계 & - & - & 47,445,440 \\\hline
\end{tabular}\par}
\medskip\noindent < 자료 2 > 감정평가 결과에 대한 피고의 이의제기 내역 \par
\medskip\noindent 1. 이의제기 요약 \par
감정평가에 의한 연간 임대료 약 1,582만원(3년 평균)은 인근의 유사한 토지들에 대한 실제 임대료 수준에 비추어 상당히 고가로서 정상적인 수준에서 크게 벗어나 감정평가 결과를 신뢰할 수 없으므로 다른 감정인에 의한 재감정평가를 신청 \par
2. 재감정 신청 증빙자료 \par
\hspace*{1em}가. 임대차계약서 사본 \par
\hspace*{2em}ㄱ) 소재지: 경기도 K시 H동 78 \par
\hspace*{2em}ㄴ) 지목 및 이용상황: 전 / 전 \par
\hspace*{2em}ㄷ) 면적: 1,998m\textsuperscript{2} \par
\hspace*{2em}ㄹ) 토지이용계획사항: 자연녹지지역, 개발제한구역 \par
\hspace*{2em}ㅁ) 주요 계약 내용 \par
{\small\centering
\begin{tabular}{|c|l|}\hline
계약기간 & 2018. 3. 10 \textasciitilde{} 2021. 3. 9 \\\hline
임대료 & 9,000,000원/연(보증금 없는 상태 기준) \\\hline
특약사항 & \parbox{0.6\linewidth}{임대차 계약기간 중 상호 합의되는 경우를 제외하고는 계약조건의 변경은 없는 것으로 하고, 임대료는 계약기간 중 매년 초일에 연간임대료 지급} \\\hline
\end{tabular}\par}
\hspace*{1em}나. 공인중개사의 사실확인서 사본: 2018년 기준 본건 및 임대차사례 토지가 소재한 지역의 연간 임대료(보증금이 없는 상태 기준)는 토지면적 약 660m\textsuperscript{2} 기준 300만원 수준이고, 최근에는 350만원 내외의 수준임을 확인함 \par
\medskip\noindent < 자료 3 > 실지조사, 자료수집 및 검토내용(요약) \par
\medskip\noindent \hspace*{1em}가. 본건 토지는 수도권 도시 근교 개발제한구역 내에 소재하며, 원예농업에 할당된 농지로 대부분 이용중임 \par
\hspace*{1em}나. 원예농업에 할당된 획지 규모는 일반적으로 약 1,600m\textsuperscript{2} \textasciitilde{} 3,300m\textsuperscript{2}임 \par
\hspace*{1em}다. 2018년 초 기준 인근 원예농업을 위한 농지 임대차에 있어 전형적인 임대차 조건은, 계약기간 3년에 연간 임대료(보증금 없는 상태 기준, 계약기간 중 매년도 초일에 연간임대료 지불)는 @4,500원/m\textsuperscript{2} 수준이었고 이후 도시화에 따른 지가 상승 등의 영향으로 지속적으로 상승하여 2021년 5월 이후에는 @5,400원/m\textsuperscript{2} 수준을 나타내고 있으며, 이와 같은 인근지역의 임대차시장 상황은 당분간 지속될 것으로 판단됨 \par
\hspace*{1em}라. 본건 및 임대차사례 토지가 속한 지역은 지속적인 도시화의 영향으로 2018년 5월 이후 2021년 4월까지 토지가격은 약 20\% 상승하였고 이러한 추세는 향후 지속될 것으로 판단됨 \par
\hspace*{1em}마. 당초 감정평가서 상의 연도별 기초가격은 시장가치에 부합하는 적정한 것으로 판단됨 \par
\hspace*{1em}바. 본건 토지 개별공시지가 \par
{\small\centering
\begin{tabular}{|c|c|}\hline
공시기준일 & 개별공시지가(원/m\textsuperscript{2}) \\\hline
2018. 1. 1. & 360,000 \\\hline
2019. 1. 1. & 382,000 \\\hline
2020. 1. 1. & 406,000 \\\hline
2021. 1. 1. & 431,000 \\\hline
\end{tabular}\par}
\hspace*{1em}사. 농지의 재산세율은 0.07\%이고 과세표준액은 시가표준액(개별공시지가)의 70\%이며, 재산세 부과 시 20\%의 지방교육세가 부가됨 \par
\hspace*{1em}아. 본건 및 임대사례토지의 개별요인은 유사하며 공히 인근지역의 일반적·평균적인 수준을 나타내고 있음 \par
\end{longquestion}
